\documentclass[11pt]{amsart}
\usepackage{amsmath,amssymb,amsthm}
\usepackage[margin=1.05in]{geometry}
\raggedbottom

\numberwithin{equation}{section}
\newtheorem{theorem}{Theorem}[section]
\newtheorem{proposition}[theorem]{Proposition}
\newtheorem{lemma}[theorem]{Lemma}
\newtheorem{corollary}[theorem]{Corollary}
\theoremstyle{definition}
\newtheorem{definition}[theorem]{Definition}
\theoremstyle{remark}
\newtheorem{remark}[theorem]{Remark}
\newtheorem{example}[theorem]{Example}

\newcommand{\Z}{\mathbb{Z}}
\newcommand{\Q}{\mathbb{Q}}
\newcommand{\R}{\mathbb{R}}
\newcommand{\N}{\mathbb{N}}
\newcommand{\C}{\mathbb{C}}
\newcommand{\cA}{\mathcal{A}}
\newcommand{\Sym}{\operatorname{Sym}}
\newcommand{\Tr}{\operatorname{tr}}
\newcommand{\Frob}{\operatorname{Frob}}
\newcommand{\Gal}{\operatorname{Gal}}
\newcommand{\Prob}{\operatorname{Prob}}
\newcommand{\ST}{\mathrm{ST}}
\newcommand{\id}{\mathrm{id}}

\PassOptionsToPackage{hyphens}{url}
\usepackage[hidelinks,bookmarks=false]{hyperref}

\title[Sato--Tate families]{Sato--Tate families in localized Bost--Connes systems:\\ open phase, no congruence obstruction,\\ and the symmetric power expansion}
\author{Ruqing Chen}
\address{GUT Geoservice Inc., Montreal, Canada}
\email{ruqing@hotmail.com}
\thanks{Sources, code and data for the entire series: \url{https://github.com/Ruqing1963/localized-bost-connes}; this paper is in \texttt{papers/XXV-sato-tate/}.}
\date{}

\begin{document}

\begin{abstract}
Let $E/\Q$ be a non-CM elliptic curve, $a_p=2\sqrt p\cos\theta_p$, and
$S_{[a,b]}=\{p:\theta_p\in[a,b]\}$ for $0\le a<b\le\pi$. We settle the three questions posed
for the localized system $\cA_{\Q,S_{[a,b]}}$.

\emph{The phase is open, and necessarily so.} By Sato--Tate $S_{[a,b]}$ has positive natural
density $\delta=\mu_\ST([a,b])=\frac1\pi\bigl[(b-a)-\tfrac12(\sin2b-\sin2a)\bigr]$ among the
primes, so $\pi_S(x)\sim\delta x/\log x$: the exponent is $\theta=1$ and the logarithmic power
is $k=1$. By Proposition 1.1 of Paper XXIV the critical series is then
$\asymp\delta\log\log x$, divergent, and the symmetry-broken phase is $(1,\infty)$, open. The
precise asymptotic is
\[
\zeta_{S_{[a,b]}}(\beta)=\bigl(\delta+o(1)\bigr)\log\frac{1}{\beta-1}\qquad(\beta\to1^+),
\]
with $o(1)$ replaced by $O(1/\log\frac{1}{\beta-1})$ under any effective Sato--Tate error
term $O(x/\log^{1+\varepsilon}x)$.
This is not an accident of the interval: \emph{a set of positive density can never have a
closed phase}, closedness requiring $k\ge2$, hence density zero with an extra logarithm.

\emph{There is no congruence obstruction.} The residue $p\bmod m$ is
$\det\bar\rho_{E,m}(\Frob_p)$, the cyclotomic character, while $\theta_p$ is governed by the
trace; Serre's open image theorem makes the mod-$m$ image open, and joint equidistribution of
$(\theta_p,p\bmod m)$ follows from the non-vanishing of $L(s,\Sym^kE\otimes\chi)$ on
$\Re s=1$, now unconditional for non-CM $E/\Q$ by Newton--Thorne together with Shahidi.
Hence every nontrivial character has a detecting set of positive density, and
\textbf{the $\mathrm{KMS}_\beta$ state is unique for all $0<\beta\le1$}. This is the
opposite verdict to Theorem 2.2(1) of Paper XXIV of the series, where the Landau and
Friedlander--Iwaniec primes are trapped in the kernel class $1\bmod4$ and hence
never unique --- and it is reached by the opposite mechanism to that theorem's part (2),
where the twin and Sophie Germain primes escape obstruction while confined to a single
\emph{generating} class. Verified numerically on $y^2=x^3+x+1$ over $6055$ primes below $6\cdot10^4$:
in each of three intervals and each modulus $m\in\{3,4,5,8\}$ every coprime class carries the
uniform proportion.

\emph{The symmetric powers control the constant, not the transition.} Expanding the indicator
in the characters $U_k(\cos\theta)=\Tr\Sym^k$, which are orthonormal for $\mu_\ST$,
\[
\zeta_{S_{[a,b]}}(s)=\delta\log\frac{1}{s-1}+\sum_{k\ge1}c_k\log L\bigl(s,\Sym^kE\bigr)+O(1),
\qquad
c_k=\frac2\pi\int_a^b\sin\bigl((k+1)t\bigr)\sin t\,dt .
\]
Only the $k=0$ term, the Riemann zeta function, has a pole at $s=1$; every $\Sym^k$ with
$k\ge1$ is automorphic with $L(1,\Sym^kE)\ne0$, so contributes $O(1)$. The higher symmetric
powers therefore fix the constant in the expansion and the shape of $\Xi_\beta$, and have no
influence whatever on $\beta_c$.

We flag one technical point that the $L^2$ expansion conceals: $\mathbf 1_{[a,b]}$ is not
smooth, so the interchange of $\sum_k$ and $\sum_p$ requires a Beurling--Selberg smoothing,
standard in equidistribution but not vacuous.
\end{abstract}

\maketitle

\section{The thermodynamic baseline}

The framework is that of \cite{Main}, cited by DOI; Papers VI and XXIV of the series,
referred to by numeral, are unpublished companions, and every fact used from them is
restated or verified in place.

\begin{proposition}[Density]\label{prop:density}
For $0\le a<b\le\pi$,
\[
\delta:=\mu_\ST\bigl([a,b]\bigr)=\frac{2}{\pi}\int_a^b\sin^2t\,dt
=\frac{1}{\pi}\Bigl[(b-a)-\tfrac12\bigl(\sin2b-\sin2a\bigr)\Bigr]>0 .
\]
By the Sato--Tate theorem for non-CM elliptic curves over $\Q$ \cite{BLGHT}, $S_{[a,b]}$ has natural
density $\delta$ in the primes, so $\pi_S(x)\sim\delta\,x/\log x$.
\end{proposition}

\begin{theorem}[Open phase]\label{thm:open}
$\beta_c(S_{[a,b]})=1$, and
\begin{equation}\label{eq:asym}
\zeta_{S_{[a,b]}}(\beta)=\sum_{p\in S_{[a,b]}}p^{-\beta}=\bigl(\delta+o(1)\bigr)\log\frac{1}{\beta-1}
\qquad(\beta\to1^+),
\end{equation}
so that $\zeta_S(1)=\infty$: the symmetry-broken phase is the \emph{open} half-line
$(1,\infty)$. If moreover $\pi_S(x)=\delta\pi(x)+O(x/\log^{1+\varepsilon}x)$ for some
$\varepsilon>0$ --- an effective Sato--Tate error term --- then the $o(1)\log\frac1{\beta-1}$
sharpens to $O(1)$.
\end{theorem}

\begin{proof}
In the notation of Proposition 1.1 of Paper XXIV, $\theta=1$ and $k=1$, so $\beta_c=\theta=1$
and the critical series is $\asymp\delta\int_2^xdt/(t\log t)=\delta\log\log x\to\infty$. For
\eqref{eq:asym}, write $E(x)=\pi_S(x)-\delta\pi(x)=o(x/\log x)$ and integrate by parts:
$\sum_{p\in S}p^{-\beta}=\delta\sum_pp^{-\beta}+\int_2^\infty x^{-\beta}\,dE(x)$, the
error being $\beta\int x^{-\beta-1}E(x)\,dx=o\bigl(\log\frac1{\beta-1}\bigr)$ from
$E=o(x/\log x)$, and $O(1)$ when $E\ll x/\log^{1+\varepsilon}x$, since
$\int x^{-2}\cdot x/\log^{1+\varepsilon}x\,dx<\infty$; finally
$\delta\sum_pp^{-\beta}=\delta\log\zeta(\beta)+O(1)=\delta\log\frac{1}{\beta-1}+O(1)$. Closedness would require
convergence at $\beta_c$, hence $k>1$ by Proposition 1.1 of Paper XXIV.
\end{proof}

\begin{corollary}[No positive-density set has a closed phase]\label{cor:noclosed}
If $S\subseteq P$ has positive natural density then $k=1$ and the low-temperature phase is
open. Closed phases require density zero, and specifically a counting function
$\asymp x/\log^kx$ with $k\ge2$ --- which by Remark 1.3 of Paper XXIV means two
simultaneous prime conditions. Sato--Tate families, being of positive density, are on the
opposite side of that divide from the Sophie Germain family of \cite[Thm.~5.5]{Main}.
\end{corollary}

\begin{example}\label{ex:densities}
$\delta=0.6090$ for $[\pi/3,2\pi/3]$, $\;0.0288$ for $[0,\pi/6]$, $\;0.0775$ for
$[2.4,\pi]$, $\;0.5$ for $[0,\pi/2]$. Empirical densities for $y^2=x^3+x+1$ over the
$6055$ good primes below $6\cdot10^4$: $0.6206$, $0.0276$, $0.0766$ respectively.
\end{example}

\section{No congruence obstruction}

\begin{lemma}[The residue is the determinant]\label{lem:det}
For $m\ge1$ and $p\nmid m N_E$, $\det\bar\rho_{E,m}(\Frob_p)=p\bmod m$, the mod-$m$
cyclotomic character. The angle $\theta_p$ is determined by
$\Tr\bar\rho_{E,\ell}(\Frob_p)=a_p$. Thus the residue class and the Sato--Tate angle are
carried by the determinant and the trace of the same representation.
\end{lemma}

\begin{theorem}[Joint equidistribution]\label{thm:joint}
Let $E/\Q$ be non-CM. For every $m\ge1$ the pairs $(\theta_p,p\bmod m)$ are equidistributed
in $[0,\pi]\times(\Z/m)^\times$ for $\mu_\ST\times\mathrm{unif}$. Consequently, for every
interval $[a,b]$ and every residue class $r$ coprime to $m$,
\[
\#\bigl\{p\le x:\theta_p\in[a,b],\ p\equiv r\ (m)\bigr\}\ \sim\ \frac{\delta}{\varphi(m)}\cdot\frac{x}{\log x}.
\]
\end{theorem}

\begin{proof}
By Lemma~\ref{lem:det} and Serre's open image theorem \cite{Serre} the mod-$m$ image is open, so the
determinant and trace directions are independent in the limit. Analytically, expanding the
indicator of $[a,b]\times\{r\}$ in the characters $U_k\otimes\chi$ of
$SU(2)\times(\Z/m)^\times$ reduces the statement to the non-vanishing and holomorphy of
$L(s,\Sym^kE\otimes\chi)$ on $\Re s=1$ for all $k\ge0$ and all Dirichlet $\chi$ mod $m$,
with a pole only for $k=0$, $\chi$ trivial. Automorphy of $\Sym^kE$ for all $k$ is
Newton--Thorne \cite{NT}; the non-vanishing on $\Re s=1$ of the automorphic $L$-functions and their
twists is Jacquet--Shalika \cite{JS} and Shahidi \cite{Shahidi}. Apply Wiener--Ikehara.
\end{proof}

\begin{theorem}[High-temperature uniqueness]\label{thm:unique}
For every non-CM $E/\Q$ and every $[a,b]$ with $\delta>0$, the system
$\cA_{\Q,S_{[a,b]}}$ has a unique $\mathrm{KMS}_\beta$ state for $0<\beta\le1$, and its
$\mathrm{KMS}_\beta$ simplex is $\Prob(G_S)$ for $\beta>1$. In particular $S_{[a,b]}$ carries
\emph{no} congruence obstruction in the sense of Theorem 2.2(1) of Paper XXIV.
\end{theorem}

\begin{proof}
By \cite[Thm.~3.9]{Main} uniqueness at $\beta$ requires
$\sum_{p\in D_\chi}p^{-\beta}=\infty$ for every nontrivial $\chi\in\widehat{G_S}$. For
$K=\Q$ such a $\chi$ is supported at finitely many places and its detecting set is the union
of the classes of $S_{[a,b]}$ outside $\ker\chi$ modulo some $m$; by
Theorem~\ref{thm:joint} that set has density $\delta(1-1/[\,\cdot\,])>0$ in the primes, so
its Dirichlet series diverges for $\beta\le1$ by Theorem~\ref{thm:open}. For $\beta>1$ the
whole of $\zeta_S(\beta)$ converges, so $\Xi_\beta=\widehat{G_S}$.
\end{proof}

\begin{remark}[Numerical verification]\label{rem:numjoint}
For $y^2=x^3+x+1$ (non-CM), the $6055$ good primes below $6\cdot10^4$:
\[
\begin{array}{lcccc}
[a,b] & m=3 & m=4 & m=5 & m=8\\\hline
[\pi/3,2\pi/3] & .500,.500 & .502,.498 & .246,.250,.254,.249 & .253,.248,.248,.250\\
{[0,\pi/6]} & .431,.569 & .449,.551 & .275,.222,.269,.234 & .228,.251,.222,.299\\
{[2.4,\pi]} & .478,.522 & .496,.504 & .259,.274,.235,.233 & .246,.265,.250,.239
\end{array}
\]
against uniform values $1/2$ and $1/4$. Every coprime class is occupied at roughly the
uniform proportion; the fluctuations in the two sparse intervals are consistent with the
sample sizes ($167$ and $464$ primes). Contrast Remark 2.3 of Paper XXIV, where the twin
and Sophie Germain families occupy the single class $2\bmod3$ --- a generating class, so no
obstruction --- and the Landau and
Friedlander--Iwaniec families the single kernel class $1\bmod4$, which does obstruct.
\end{remark}

\begin{remark}[Why this family passes the test]\label{rem:whypass}
The obstructed families of Paper XXIV --- Landau and Friedlander--Iwaniec --- are defined by
a condition that \emph{involves}
the residue, $p$ being a value of a polynomial, so the condition and
the residue are not independent, and the dependence lands in the kernel. (The twin and
Sophie Germain conditions are also residue-entangled, but land in a generating class and
escape.) Here the defining condition involves $a_p$, and by
Lemma~\ref{lem:det} the residue is the determinant while $a_p$ is the trace of the same
Galois representation. Open image is exactly the statement that these two coordinates are
independent. The distinction is structural, not a matter of the density being larger.
\end{remark}

\section{The symmetric power expansion}

\begin{proposition}[Chebyshev coefficients]\label{prop:cheb}
The functions $U_k(\cos t)=\sin\bigl((k+1)t\bigr)/\sin t$, $k\ge0$, are orthonormal in
$L^2(\mu_\ST)$ and satisfy $U_k(\cos\theta_p)=\Tr\Sym^k(U_p)$ where
$U_p=p^{-1/2}\rho_E(\Frob_p)$. Writing
$\mathbf 1_{[a,b]}=\sum_{k\ge0}c_kU_k$ in $L^2(\mu_\ST)$,
\[
c_0=\delta,\qquad
c_k=\frac{2}{\pi}\int_a^b\sin\bigl((k+1)t\bigr)\sin t\,dt
=\frac1\pi\Bigl[\frac{\sin kt}{k}-\frac{\sin(k+2)t}{k+2}\Bigr]_a^b\quad(k\ge1).
\]
\end{proposition}

\begin{theorem}[The expansion]\label{thm:expansion}
For $\Re s>1$,
\begin{equation}\label{eq:expansion}
\zeta_{S_{[a,b]}}(s)=\delta\,\log\frac{1}{s-1}+\sum_{k\ge1}c_k\log L\bigl(s,\Sym^kE\bigr)+O(1),
\end{equation}
the series over $k$ summed by truncation as in Remark~\ref{rem:smoothing}, and every term
with $k\ge1$ is $O(1)$ as $s\to1^+$, since $\Sym^kE$ is automorphic
(Newton--Thorne \cite{NT}) and $L(1,\Sym^kE)\ne0$ (Shahidi \cite{Shahidi}). Hence the pole, the critical exponent
$\beta_c=1$ and the divergence at $\beta=1$ all come from the $k=0$ term, the Riemann zeta
function; the symmetric powers of higher order determine the constant and the fine structure
of $\Xi_\beta$ but not the transition.
\end{theorem}

\begin{proof}
$\log L(s,\Sym^kE)=\sum_p\Tr\Sym^k(U_p)p^{-s}+O(1)$ for $\Re s>1/2$, the prime-power terms
converging by Deligne's bound. Substitute the expansion of the indicator and use
$\log\zeta(s)=\log\frac{1}{s-1}+O(1)$ for the $k=0$ term.
\end{proof}

\begin{remark}[The smoothing is not vacuous]\label{rem:smoothing}
$\mathbf 1_{[a,b]}$ is not continuous, so the Chebyshev series converges in $L^2(\mu_\ST)$
but not uniformly, and the interchange of $\sum_k$ with $\sum_p$ in \eqref{eq:expansion} is
not justified term by term. The standard remedy is to bracket the indicator between
Beurling--Selberg majorants and minorants of bounded degree, which introduces an error
$O(1/K)$ from truncating at $k\le K$ and is then optimized against the error term in the
Sato--Tate count. We note that the truncation must also absorb the $k$-dependence of the
individual $O(1)$'s: both the prime-power tails ($|\mathrm{tr}\,\Sym^k|\le k+1$) and the
size of $\log L(s,\Sym^kE)$ near $s=1$ grow with $k$ through the analytic conductor, so the
constants are not uniform and the series over $k$ converges only through the bracketing, not
term by term. The formula is correct; the proof of the interchange is the usual one and
is not automatic.
\end{remark}

\begin{example}[Coefficients]\label{ex:coeffs}
For $[a,b]=[\pi/3,2\pi/3]$: $c_0=0.608998$, $c_1=0$, $c_2=-0.413497$, $c_3=0$,
$c_4=0.137832$, $c_5=0$, $c_6=0.068916$. The odd coefficients vanish by the symmetry of the
interval about $\pi/2$; $\Sym^2E$ carries the leading correction, which is the familiar
statement that the second moment of $a_p$ is the first place a symmetric interval is
sensitive to the curve.
\end{example}

\section{Assessment}

\[
\begin{array}{lll}
\text{Question} & \text{Answer} & \text{Reference}\\\hline
\beta_c(S_{[a,b]}) & 1 & \text{Thm.~1.2}\\
\text{phase at }\beta=1 & \textbf{OPEN }(1,\infty) & \text{Thm.~1.2}\\
\zeta_S(\beta),\ \beta\to1^+ & \delta\log\frac{1}{\beta-1}+O(1) & \text{Thm.~1.2}\\
\text{positive density}\Rightarrow\text{closed?} & \textbf{never} & \text{Cor.~1.3}\\
\text{joint equidistribution} & \textbf{holds, unconditionally} & \text{Thm.~2.2}\\
\text{congruence obstruction} & \textbf{none} & \text{Thm.~2.3, Rem.~2.4}\\
\text{unique }\mathrm{KMS}_\beta\text{ for }\beta\le1 & \textbf{yes} & \text{Thm.~2.3}\\
\Sym^k\ \text{controls }\beta_c & \textbf{no; only the constant} & \text{Thm.~3.2}\\
\text{term-by-term interchange} & \text{needs smoothing} & \text{Rem.~3.3}
\end{array}
\]

No no-go theorem is warranted here, and that is itself the finding. The Sato--Tate
intervals pass the congruence test of Paper XXIV by \emph{equidistribution} --- the clean
mechanism --- where the twin and Sophie Germain families pass it, by that paper's revised
Theorem 2.2(2), through the opposite extreme of confinement to a generating class. The
reason here is structural: the defining condition is a
constraint on the \emph{trace} of the Galois representation, while the residue class is its
\emph{determinant}, and Serre's open image theorem \cite{Serre} is precisely the independence of the two.
The obstructed families of Paper XXIV --- Landau and Friedlander--Iwaniec --- fail because
their defining conditions are congruence conditions in disguise, landing in the kernel.

The price is that the phase is open. Positive density forces $k=1$ and hence divergence at the
critical point, so these systems cannot exhibit the closed low-temperature phase that
\cite[Thm.~5.5]{Main} found for the Sophie Germain family. By Corollary~\ref{cor:noclosed} the
two desiderata --- a closed phase and an unobstructed high-temperature phase --- are in
tension: the first needs density zero with two prime conditions, the second is easiest for
sets defined by trace conditions, which have positive density. By the revised Theorem 2.2(2) of Paper XXIV the twin, Sophie Germain and Chen families all
combine a closed phase with the absence of obstruction, Chen unconditionally on the
divergence side; those remain the only such families we know.

Two questions. Can one intersect a Sato--Tate interval with a sparse condition to obtain
$k=2$ while keeping joint equidistribution, for instance
$\{p:\theta_p\in[a,b],\ p+2\in P_2\}$ --- the Chen condition contributing the extra logarithm
and the trace condition remaining independent of it? And second, the anomalous sets of
Theorem 5.3 of Paper VI were built by letting $[a,b]$ shrink with $p$; by
Theorem~\ref{thm:open} any \emph{fixed} interval gives $\beta_c=1$ with an open phase, so the
whole of the phase-structure interest in Sato--Tate families lies in shrinking intervals,
which is where Paper VI placed it.

\begin{thebibliography}{9}
\bibitem{BLGHT} T. Barnet-Lamb, D. Geraghty, M. Harris, R. Taylor, \emph{A family of Calabi--Yau varieties and potential automorphy II}, Publ. RIMS 47 (2011), 29--98.
\bibitem{NT} J. Newton, J. A. Thorne, \emph{Symmetric power functoriality for holomorphic modular forms}, Publ. Math. IHES 134 (2021), 1--116.
\bibitem{Serre} J.-P. Serre, \emph{Propri\'et\'es galoisiennes des points d'ordre fini des courbes elliptiques}, Invent. Math. 15 (1972), 259--331.
\bibitem{JS} H. Jacquet, J. A. Shalika, \emph{A non-vanishing theorem for zeta functions of $GL_n$}, Invent. Math. 38 (1976/77), 1--16.
\bibitem{Shahidi} F. Shahidi, \emph{On certain $L$-functions}, Amer. J. Math. 103 (1981), 297--355.
\bibitem{Main} R. Chen, \emph{KMS state uniqueness and phase transitions in localized Bost--Connes systems}, preprint, 2026, \newline\url{https://doi.org/10.5281/zenodo.22152101}.
\end{thebibliography}

\end{document}
